% Generated draft: requires mathematical and source review.
\documentclass[11pt]{article}
\usepackage[margin=25mm]{geometry}
\usepackage{fontspec}
\setmainfont{Latin Modern Roman}
\usepackage{amsmath,amssymb,mathtools,unicode-math}
\setmathfont{Latin Modern Math}
\usepackage{enumitem,needspace,xurl,hyperref}
\setlength{\parindent}{0pt}
\setlength{\parskip}{5pt}
\setlength{\emergencystretch}{3em}
\newcommand{\N}{\mathbb{N}}
\newcommand{\Z}{\mathbb{Z}}
\newcommand{\Q}{\mathbb{Q}}
\newcommand{\R}{\mathbb{R}}
\newcommand{\C}{\mathbb{C}}
\newcommand{\F}{\mathbb{F}}
\newcommand{\A}{\mathbb{A}}
\newcommand{\PP}{\mathbb{P}}
\newcommand{\E}{\mathbb{E}}
\newcommand{\eps}{\varepsilon}
\newcommand{\dd}{\,\mathrm{d}}
\begin{document}
\section*{UnsolvedMath ID: 30004780}
\textbf{Source problem code:} catalogue entry
\subsection{Definitions and mathematical statement}
Let \(F\) be a local or global field, and let \(E\) and \(K\) be quadratic étale \(F\)-algebras. Write \(a\mapsto\bar a\) for the canonical involution of \(E/F\). Under the usual interpretation of a skew-Hermitian space, \(V\) is a finite-dimensional \(E\)-module with a nondegenerate sesquilinear form \(h\) satisfying
\[
 h(av,bw)=a\,h(v,w)\,\bar b,
 \qquad h(w,v)=-\overline{h(v,w)}.
\]
Its unitary group consists of the \(E\)-linear automorphisms preserving \(h\). Interpret
\[
 V_K=V\otimes_F K
\]
as a skew-Hermitian space over \(E\otimes_F K\), with the scalar-extended form and involution \(\bar{\phantom{a}}\otimes\operatorname{id}_K\). Scalar extension gives the natural inclusion of \(U(V)\) into the group denoted \(U(V_K)\) in the source. These are conventional interpretations of notation that the supplied record does not explicitly define.

For local \(F\), let \(\Pi\) be an irreducible representation of \(U(V_K)\), and let \(\omega_{V,\psi,\mu}\) denote the Weil representation of \(U(V)\) specified in the source. Once an appropriate representation category and the auxiliary data \(\psi,\mu\) have been fixed, define the branching space and its multiplicity by
\[
 \mathcal H_V(\Pi;\psi,\mu)
 =\operatorname{Hom}_{U(V)}
 \bigl(\Pi|_{U(V)},\omega_{V,\psi,\mu}\bigr),
 \qquad
 m_V(\Pi;\psi,\mu)
 =\dim_{\C}\mathcal H_V(\Pi;\psi,\mu).
\]
Here the Hom space consists of intertwining maps in the chosen category; the record does not specify that category or establish finiteness of its dimension.

\textbf{Requested local and global assertions.} For the data to which the twisted Fourier--Jacobi Gan--Gross--Prasad conjectures apply, determine \(m_V(\Pi;\psi,\mu)\) locally and determine the corresponding global period when \(F\) is global, in accordance with the conjectural twisted GGP formulas. The supplied record does not reproduce those formulas or define the global period. It also omits the restrictions on the fields, representations, and auxiliary characters needed for their precise formulation. Consequently, this is a statement of the supplied research target, not a complete conjectural identity: neither a specific multiplicity formula nor a specific global period identity can be recovered from this record alone.
\subsection{Short English statement}
The problem asks for the twisted Fourier--Jacobi Gan--Gross--Prasad formulas describing the restriction of \(\Pi\) from \(U(V_K)\) to \(U(V)\), measured against a Weil representation, and the corresponding global period. The supplied record omits the formulas and several necessary hypotheses, so the exact conjectures require human mathematical review.
\subsection{Sources}
\begin{itemize}
\item Supplied local UnsolvedMath registry record, permanent ID \(30004780\), title: \emph{The Twisted Fourier--Jacobi Gan--Gross--Prasad Conjecture}; exact source problem code: \texttt{catalogue entry}. The recorded proposed year is 2021; no individual proposer is supplied.
\item Supplied report citation: \emph{Automorphic Forms, Geometry and Arithmetic (2022)}, Oberwolfach Reports, pp.\ 2089--2156, DOI \texttt{10.4171/owr/2021/39}. \url{https://doi.org/10.4171/owr/2021/39}. The background describes Oberwolfach Reports as documenting workshops and their open-problem discussions.
\item Supplied literature reference: \emph{Twisted GGP Problems and Conjectures}. \url{https://arxiv.org/abs/2204.10108}.
\item Supplied literature reference: \emph{Local twisted GGP for U(V\_K)/U(V)}. \url{https://arxiv.org/abs/2511.01301}.
\item Supplied literature reference: \emph{Twisted GGP for unramified quadratic extensions}. \url{https://arxiv.org/abs/2307.15234}.
\item Supplied literature assessment, checked \texttt{2026-08-22}: the record labels the problem partially solved and reports precise formulations and verified low-rank and principal-series evidence, while describing the general local and global conjectures as open. This is supplied dated curation derived from the OWR record, not an independently verified literature assessment; its time-sensitive status and cited claims require independent verification.
\item The registry labels the statement \texttt{corrected\_verified} and reports that its standalone formulation and hypotheses were checked against the official Oberwolfach report. This verification claim is supplied metadata and has not been independently checked here.
\end{itemize}
\end{document}
